\documentclass[12pt]{article}

\usepackage[a4paper,margin=2.5cm]{geometry}
\usepackage{fontspec}
\usepackage{polyglossia}
\usepackage{amsmath,amssymb}
\usepackage{enumitem}

\setdefaultlanguage{russian}
\setmainfont{Times New Roman}
\setsansfont{Arial}
\setmonofont{Menlo}
\newfontfamily\cyrillicfont{Times New Roman}
\newfontfamily\cyrillicfontsf{Arial}
\setlist[itemize]{leftmargin=2em}

\begin{document}

\begin{center}
    {\LARGE \textbf{Смысл использования моделей волатильности и ценообразования опционов}}
\end{center}

\section*{Предположение исследования}

В работе принимается базовое допущение:

\begin{itemize}
    \item рыночная цена опциона \(P^{market}\) рассматривается как наблюдаемый эталон (benchmark);
    \item модельная цена \(P^{model}\) --- как приближение этой величины в рамках заданной модели;
    \item качество модели тем выше, чем меньше её отклонение от рыночной цены.
\end{itemize}

Формально, для двух моделей \(M_1\) и \(M_2\) выполняется критерий предпочтения:

\[
Error_{M_1} < Error_{M_2} \Rightarrow M_1 \text{ описывает рынок лучше}
\]

где ошибка может измеряться, например, как:

\[
Error^{price} = |P^{model} - P^{market}|
\]

Данное предположение задаёт единый критерий сравнения моделей и используется во всех последующих разделах.

\section*{1. Проблема: есть цена, но нет её смысла}

На рынке опционов наблюдаются реальные цены \(P^{market}\), однако сами по себе они не являются интерпретируемыми.

Цена опциона агрегирует множество факторов:

\begin{itemize}
    \item ожидания волатильности;
    \item временную стоимость;
    \item вероятности будущих движений;
    \item рыночные эффекты (ликвидность, спрос и предложение).
\end{itemize}

Без модели невозможно ответить на ключевые вопросы:

\begin{itemize}
    \item высокая ли эта цена или нет;
    \item за счёт чего она сформирована;
    \item как она изменится при изменении параметров рынка.
\end{itemize}

\section*{2. Роль модели: структурирование цены}

Модель ценообразования задаёт функциональную зависимость:

\[
P^{model} = f(S, K, T, \sigma, r)
\]

где:

\begin{itemize}
    \item \(S\) --- цена базового акива;
    \item \(K\) --- страйк;
    \item \(T\) --- время до экспирации;
    \item \(\sigma\) --- волатильность;
    \item \(r\) --- безрисковая ставка.
\end{itemize}

Модель позволяет:

\begin{itemize}
    \item перейти от цены к подразумеваемой волатильности \(\sigma^{impl}\);
    \item анализировать чувствительности (Greeks);
    \item интерпретировать цену через параметры.
\end{itemize}

Важно: модель не <<разбивает>> цену на части, а задаёт \textbf{структуру зависимости}.

\section*{3. Позиция исследования}

В рамках работы принимается следующая установка:

\begin{itemize}
    \item рыночная цена \(P^{market}\) --- наблюдаемый эталон;
    \item модельная цена \(P^{model}\) --- приближение.
\end{itemize}

Модель не рассматривается как источник <<истинной>> цены, превосходящей рыночную.

Её задача --- \textbf{объяснить наблюдаемую цену в рамках своих предположений}.

\section*{4. Цель анализа моделей}

Основной вопрос:

\begin{quote}
Насколько хорошо модель описывает наблюдаемые рыночные цены и сохраняет ли это качество в разных рыночных условиях?
\end{quote}

\bigskip

\begin{center}
    {\LARGE \textbf{Оценка качества моделей через ошибку и её устойчивость}}
\end{center}

\section*{1. Метрика качества: ошибка модели}

Качество модели оценивается через её отклонение от рынка.

\subsection*{Основная метрика}

\[
Error^{price} = |P^{model} - P^{market}|
\]

\subsection*{Альтернативная метрика (через волатильность)}

\[
Error^{IV} = |\sigma^{model} - \sigma^{market}|
\]

где \(\sigma^{market}\) --- implied volatility, извлечённая из рыночной цены.

\section*{2. Что именно анализируется}

Задача исследования --- не просто измерить ошибку, а изучить зависимость:

\[
Error = f(\text{market regime})
\]

\section*{3. Понятие устойчивости}

Под устойчивостью модели понимается способность модели сохранять низкую ошибку при изменении рыночных условий.

\section*{4. Роль рыночных режимов}

Рассматриваются различные режимы:

\begin{itemize}
    \item низкая волатильность (спокойный рынок);
    \item высокая волатильность (стресс, кризис).
\end{itemize}

Ожидаемое поведение:

\begin{itemize}
    \item в спокойных условиях:
\end{itemize}

\[
Error \text{ --- низкая}
\]

\begin{itemize}
    \item в стрессовых условиях:
\end{itemize}

\[
Error \text{ --- возрастает}
\]

Причина: нарушаются предположения моделей (например, constant volatility в Black--Scholes).

\section*{5. Сравнение моделей}

Разные модели демонстрируют различную устойчивость:

\begin{itemize}
    \item простые модели (Black--Scholes):
    \begin{itemize}
        \item хорошо работают в стабильных условиях;
        \item теряют точность в кризисах;
    \end{itemize}
    \item более сложные модели (Heston, regime-switching):
    \begin{itemize}
        \item лучше учитывают изменчивость рынка;
        \item потенциально дают меньшую ошибку в стрессовых режимах.
    \end{itemize}
\end{itemize}

\section*{6. Итоговая постановка задачи}

Центральная задача исследования:

\[
\text{сравнить модели по функции } Error = f(\text{regime})
\]

и определить:

\begin{itemize}
    \item какие модели лучше аппроксимируют рынок;
    \item какие модели более устойчивы к смене рыночных режимов.
\end{itemize}

\section*{7. Эконометрические подходы к анализу ошибок}

Для формального анализа качества моделей можно использовать простые эконометрические методы без сложной математики.

\subsection*{7.1. Регрессионный анализ}

Идея: проверить, от каких факторов зависит ошибка модели (например, от уровня implied volatility, реализованной волатильности и рыночного режима).

Это позволяет:

\begin{itemize}
    \item понять, какие факторы увеличивают ошибку;
    \item проверить, значимы ли эти факторы;
    \item описать, в каких условиях модель работает хуже.
\end{itemize}

\subsection*{7.2. Анализ по нескольким опционам (panel data)}

Если есть данные по разным страйкам и датам, можно анализировать ошибку одновременно по множеству опционов.

Это позволяет:

\begin{itemize}
    \item учитывать различия между опционами;
    \item отделить эффекты, связанные с конкретным опционом, от общерыночных;
    \item получить более устойчивые выводы.
\end{itemize}

\subsection*{7.3. Анализ динамики ошибки во времени}

Рассматривается, как ошибка меняется со временем.

Это позволяет:

\begin{itemize}
    \item увидеть, есть ли периоды, когда ошибка резко возрастает;
    \item проверить, усиливается ли ошибка в стрессовые периоды;
    \item выявить кластеризацию ошибок (периоды, когда они последовательно высокие).
\end{itemize}

\subsection*{7.4. Сравнение рыночных режимов}

Простейший подход --- разделить выборку на два режима:

\begin{itemize}
    \item спокойный рынок;
    \item стрессовый рынок.
\end{itemize}

Далее сравнивается средний уровень ошибки в этих режимах.

Ключевая идея:

\begin{itemize}
    \item если ошибка в стрессовом режиме выше, модель неустойчива;
    \item если ошибка стабильна, модель более устойчива.
\end{itemize}

\subsection*{7.5. Итоговые метрики качества}

Для агрегированной оценки используются простые показатели:

\begin{itemize}
    \item средняя абсолютная ошибка (MAE);
    \item средняя квадратичная ошибка (RMSE);
    \item ошибка в implied volatility.
\end{itemize}

Эти метрики позволяют сравнивать модели между собой и делать итоговые выводы.

\section*{Итог}

Модель оценивается не по прибыли, а по тому, \textbf{насколько точно и стабильно она объясняет рыночные цены опционов}.

\end{document}
